% =============================================================================
% Draft mục 3.3.7 — Fallback Bridge Aggregator (DESIGN-ONLY)
% =============================================================================

\subsection{Fallback Bridge Aggregator}
\label{sec:fallback-bridge}

\noindent\textbf{Lưu ý phạm vi.} Mục này trình bày
\emph{thiết kế ý đồ} của tầng fallback. Trong scope của luận văn, cài
đặt thực tế chưa được hoàn tất; mã nguồn hiện chỉ chứa interface và
một adapter mẫu. Phần thực thi đầy đủ và đo benchmark được liệt kê
trong Chương Future Work.

\subsubsection*{Lý do tồn tại}

Matching Engine (\S\ref{sec:matching-engine}) chỉ thanh toán được
phần intent có thể tham gia vào một chu trình thoả ngưỡng
$\mathrm{ratio} > x$. Có hai nhóm trường hợp một intent
\emph{không thể} được khớp bởi tầng netting:

\begin{enumerate}
  \item \textbf{Không có đối ứng.} Tại thời điểm intent được đưa
    vào orderbook, không có intent nào trên cặp chuỗi đối ứng có khả
    năng tạo cycle với nó.
  \item \textbf{Cân bằng lệch.} Cycle tồn tại nhưng
    $\mathrm{ratio}$ quá nhỏ — phần khớp được không bù được chi phí
    gas, hoặc phần dư còn lại quá lớn so với nhu cầu của người dùng.
\end{enumerate}

Trong các trường hợp đó, người dùng vẫn cần một con đường để hoàn
tất việc bridge. \emph{Fallback Bridge Aggregator} là tầng cuối
cùng: thay vì cố matching P2P, nó định tuyến intent qua một bridge
"thông thường" — một bridge bên thứ ba đã có thanh khoản thực sự
trên cả hai chuỗi — và chấp nhận trả phí cao hơn để đổi lấy
\emph{tính chắc chắn về thời gian hoàn tất}.

\subsubsection*{Vị trí trong mô hình thực thi ba lớp}

Theo mô hình thực thi ba lớp (\S "Mô hình thực thi ba lớp"), một
intent điển hình đi qua các lớp theo thứ tự ưu tiên:

\begin{enumerate}
  \item \textbf{Netting Layer.} Đợi tối đa $T_{\text{net}}$ giây (mặc
    định: vài tick scheduler). Nếu được matching engine khớp một
    phần hoặc toàn bộ, phần đã khớp đi qua Execution Engine và HTLC.
    Phần dư (nếu có) tiếp tục.
  \item \textbf{Intent Accelerator Layer} (chưa nằm trong scope). Mời
    các solver bên ngoài quote một mức giá riêng cho phần dư;
    người gửi có thể chấp nhận hoặc từ chối. Nếu chấp nhận và solver
    thực hiện, phần dư được thanh toán qua HTLC như một flow đặc
    biệt.
  \item \textbf{Fallback Bridge Aggregator.} Nếu cả hai lớp trên
    không xử lý được trong thời gian $T_{\text{total}}$, phần intent
    còn lại được định tuyến qua một bridge bên thứ ba.
\end{enumerate}

Trật tự này đảm bảo: phần lớn intent hưởng lợi từ phí thấp + thanh
toán P2P (lớp 1), một số ít chấp nhận trade-off (lớp 2), và không
có intent nào bị "treo" vô thời hạn (lớp 3).

\subsubsection*{Các bridge ứng cử viên}

Thiết kế tham khảo ba bridge phổ biến để tích hợp:

\begin{itemize}
  \item \textbf{Across Protocol.} Mô hình intents-based, dùng các
    relayer cạnh tranh. Phù hợp với các cặp Ethereum L1 $\leftrightarrow$ L2 (đặc
    biệt là Arbitrum, Base, Optimism). Phí thường thấp; thời gian
    finalize tính bằng phút.
  \item \textbf{LayerZero.} Generic messaging layer. Tích hợp tốt với
    nhiều chuỗi non-EVM trong tương lai. Mô hình endpoint + DVN
    (Decentralized Verifier Network), tuỳ vào lựa chọn DVN mà chi
    phí và độ trễ khác nhau.
  \item \textbf{Chainlink CCIP.} Cross-Chain Interoperability
    Protocol. Phù hợp khi yêu cầu finality cao và sẵn sàng trả phí
    đáng kể. Mạng đã được tích hợp trên cả ba chuỗi mục tiêu.
\end{itemize}

Tầng aggregator được thiết kế \emph{không khoá vào một bridge cụ
thể}; thay vào đó, nó định nghĩa một interface chung và mỗi bridge
là một adapter implement interface đó.

\subsubsection*{Interface đề xuất}

Theo nguyên tắc đa hình ngầm (structural typing) của TypeScript,
interface tối thiểu của một bridge adapter có dạng:

\begin{verbatim}
interface BridgeAdapter {
  name: string;

  // Lấy quote: phí, ETA, lãi suất bảo đảm
  quote(req: QuoteRequest): Promise<QuoteResponse>;

  // Phát giao dịch bridge; trả về txHash trên chuỗi nguồn
  // và thông tin để theo dõi đến khi hoàn tất trên chuỗi đích
  send(req: SendRequest): Promise<SendResult>;

  // Theo dõi trạng thái của một bridge đang chạy
  status(handle: BridgeHandle): Promise<BridgeStatus>;
}
\end{verbatim}

\texttt{Aggregator} duy trì danh sách các \texttt{BridgeAdapter}, gọi
\texttt{quote} song song, chọn quote tốt nhất theo một hàm chi phí
(thường là $\mathrm{fee} + \mathrm{ETA} \cdot \alpha$ cho hằng số
$\alpha$ phản ánh "giá trị thời gian" của người dùng), rồi gọi
\texttt{send} trên
adapter tương ứng. Trạng thái sau đó được theo dõi và cập nhật vào
\texttt{ExecutionRecord} qua cùng kênh sự kiện như nhánh netting.

\subsubsection*{Câu hỏi mở (đặt cho Future Work)}

\begin{itemize}
  \item \textbf{Tiêu chí kích hoạt fallback.} Hiện thiết kế dựa vào
    timeout đơn lẻ; có thể nâng cấp thành chính sách động dựa trên
    profile rủi ro của người dùng (ví dụ: dùng fallback ngay nếu
    intent có \texttt{deadline} quá gần).
  \item \textbf{Lỗ hổng khi quote.} Quote đáo hạn nhanh; aggregator
    cần xử lý trường hợp quote hết hiệu lực giữa lúc người dùng
    chấp nhận và lúc \texttt{send}. Phương án đề xuất: re-quote và
    yêu cầu re-confirm nếu lệch quá ngưỡng.
  \item \textbf{Tích hợp với accounting nội bộ.} Vì fallback tốn
    phí thực, hệ thống cần ghi rõ phần phí được chuyển cho bridge
    bên thứ ba để đối chiếu báo cáo P\&L. Yêu cầu một bảng phụ
    \texttt{fallback\_executions} trong lược đồ dữ liệu.
\end{itemize}
